\section{Literature Review}
\label{sec:lit}

% Rebuilt 12 Oct 2026 to the thesis plan of 11 Oct (section 4, "2 Literature"): 2.1 transport and link models,
% 2.2 membership change with the numeric table, 2.3 capture proof and docking, 2.4 switching and runtime assurance,
% 2.5 the execution layer, 2.6 simulation to reality, 2.7 the gap. modelling.tex and control.tex are no longer
% input: their link-model and centralised text is condensed into lit_transport.tex, their membership and docking
% text into lit_membership.tex and lit_capture.tex, and the distributed detail with the architecture figure is
% Appendix F (app_network_lit.tex). The old comparison table (tab:control_comparison) is replaced by the
% membership-change table (tab:membership). Template note: develop from the previous reports and the supervisors'
% feedback; flow logically, not necessarily chronologically; end with the major work, its relevance, and the gaps.

This section reviews how a team carries a suspended payload and models its links
(Section~\ref{sec:lit_models}), what is known about a vehicle leaving or joining a carrying team
(Section~\ref{sec:lit_team}), how a catch is proven and vehicles dock in flight (Section~\ref{sec:lit_capture}), how
control is switched safely (Section~\ref{sec:lit_switching}), the execution layer on an unmodified flight controller
(Section~\ref{sec:lit_execution}) and what a simulator can be trusted to predict (Section~\ref{sec:lit_sim}), and ends
with the gap (Section~\ref{sec:lit_gap}).

% lit_transport.tex -- Literature 2.1 (v3_50: condensed from v3_full).

\subsection{Cooperative transport and link models}
\label{sec:lit_models}

Several quadcopters carrying one suspended payload raise the mass a team can lift and can control its full pose, at
the price of a coupled system in which every vehicle feels the others through the load \cite{tudelft,
mohammadiDecentralizedMotionControl2018}. Most of the field, reviewed by Lv et al.\ \cite{lv2026review} and Estevez
et al.\ \cite{estevezReviewAerialTransportation2024}, concerns tracking, tension allocation and disturbance rejection
for a team whose membership never changes, validated mostly in simulation.

\paragraph{Link models.} The common choice is the rigid massless rod: a link of fixed length along a unit vector,
carrying a tension $\lambda_i \ge 0$, so that it can pull but not push \cite{efficient, wahbaKinodynamicMotionPlanning2024}.
It is exact for a taut cable under a hovering vehicle \cite{michael2011cooperative}, and the non-negativity of the
tension is a modelling choice about when the link is believed \cite{six2018flying}. Mass-spring-damper chains capture
sag at six states per element \cite{distributedNMPC}; a slack link, a unilateral constraint, is still open
\cite{efficient}. The rig of this thesis carries its ring on rigid rods, so the rod model is the hardware's own, except
for the newcomer's tether (Section~\ref{sec:overview}).

\paragraph{Centralised and distributed control.} A central node with the global state can impose tension limits,
collision avoidance and a force allocation directly, at the cost of a solve that grows with the team and a single point
of failure \cite{nathanReviewAutonomousMultiagent2011, dimakosStudyCentralisedDecentralised2021}. The benchmark for
this thesis is Sun et al.\ \cite{tudelft}, who solve a centralised optimal control problem over the full
load--cable--vehicle dynamics at 10\,Hz and track the plans onboard with incremental nonlinear dynamic inversion (INDI),
which compensates the cable force from inertial measurements; three and four vehicles reach 5\,m/s. Two things are left
open there that matter here: the solve time grows exponentially with the team, and the results do not separate what the
INDI force compensation contributes from what the planner does. Full-pose MPC scales to ten vehicles in simulation
\cite{sunNonlinearMPCFullPose2023}, and allocation can be separated as quadratic programs
\cite{efficient, vazquez-soquiSQPBasedCableTensionAllocation2026}; all of these hold the team fixed. Without a central
node each vehicle solves for itself from what it measures or what its neighbours send
\cite{scattoliniArchitecturesDistributedHierarchical2009, distributedNMPC, tasoojiEventTriggeredNonlinearModel}. Closest
in spirit to a changing team is the passivity-based control of Quan et al.\ \cite{quan2025swarm, liu2026passivity}, in
which each vehicle damps its motion relative to the payload from local measurements, so that a member leaving injects
energy that the damping removes. The project plan asked whether such a network or the trajectory-based planner better
balances robustness and tracking; that was answered on an earlier plant and is recorded in
Appendix~\ref{sec:app_network}. This thesis takes the optimal control problem of Sun et al.\ as its planner, with the
team size and the attachment geometry made runtime parameters (Section~\ref{sec:model}), and a tracking controller per
vehicle on the ground in place of INDI, because the flight controllers are unmodified (Section~\ref{sec:lit_execution}).


% lit_membership.tex -- Literature 2.2: changes of team membership, with the numeric table T3 (drafted 12 Oct 2026
% from the thesis plan of 11 Oct). The facts on Quan et al. and Zeng et al. were verified from the full texts on
% 11 Oct (hand-back docs/thesis_work_2026-10-11.md); the cells in brackets wait for a full-text read (plan C2-A:
% Bacheti and Delbene were read from abstracts only) or for the October flights. Entries the plan lists for this
% section that are not yet in the bibliography (Oishi 2023, Gao 2024, Delbene 2024/2025, Bacheti 2026, Xu 2025) wait for the approval batch;
% Fink 2011, Tognon 2018 and Farina 2020 were added 12 Oct.
% Papers found with Undermind and approved by Wesley, 10 Oct (DOIs, abstracts and sources in
% docs/thesis_drafts/refs_2026-10-10.md; alternates listed there).

\subsection{Changes of team membership}
\label{sec:lit_team}

In this thesis a reconfiguration is a change of team membership: a vehicle leaving or joining while the payload
is carried. The word is also used for a fixed team changing its shape, as when a team contracts to pass a doorway
\cite{jackson2020scalable, tensionbased}; that is a formation change, and the team's force allocation keeps its
dimension. A membership change does not. Most of the work of Section~\ref{sec:lit_models} assumes a fixed team
\cite{efficient, wahbaKinodynamicMotionPlanning2024, quan2025swarm}; the few works that change it are set out in
Table~\ref{tab:membership}, with the detection of the change kept apart from the recovery from it, because the
two are often reported as one time.

\paragraph{A vehicle leaving, on hardware.} Quan et al.\ \cite{quan2025swarm} remove a vehicle from a carried
load in flight: four robots carry a 4\,kg load, one releases its cable on command, and the other three
reorganise to a near-equilateral layout and rebalance their tensions within 13\,s, decentralised and without
replanning. They report the cable tensions and one robot's velocity through the event, not the load's attitude.
Zeng et al.\ \cite{zeng2025marl} switch one of three vehicles off in flight and leave it hanging from a 1.4\,kg
load; the other two keep five degrees of freedom of the load under a learned decentralised policy that observes
only the load's pose, and the load's position and attitude are plotted through the loss with no transient
number. Bacheti et al.\ \cite{tensionbased} handle a removal by repositioning the remaining vehicles before the
release so that the transient is small, on hardware with a point-mass payload; [whether the agent count changes
in their experiments is to be verified from the full text]. So on hardware a vehicle has been released from a
carried load and a vehicle has been lost from one, and in neither case is the load's attitude through the event
given as a number.

\paragraph{A link that fails, in simulation.} A cable that breaks rather than releases is treated as a fault: the
remaining tensions are estimated and the load reassigned \cite{liang2021fault}, a supervisor reconfigures a
passivity-based formation \cite{saiella2021fault}, or each vehicle's measured tension is routed into its thrust so
that five vehicles recover from one or two severed cables within a pendulum period \cite{hajieghrary2026passive}; see
also \cite{farrasFaultTolerantControlMultiQuadcopter2024}. None flies, where a time is given it is the recovery and not
the detection, and none adds a vehicle.

\paragraph{A vehicle joining.} No work attaches a vehicle to a payload that a team is already carrying. Quan et al.\
describe a joining robot only in prose \cite{quan2025swarm}, and the event-triggered architecture of Tasooji et al.\
\cite{tasoojiCooperativeControlMultiQuadrotors} is argued to scale but not analysed for a varying team. The closest
hardware precedent is a payload handed from one multirotor to another, the grasp confirmed by a magnetic sensor before
the giver releases \cite{shankar2020inflight}; the team does not grow and the payload is not shared.

\paragraph{Prepared changes and statics.} Plug-and-play MPC lets subsystems join or leave a running network without
losing stability, provided each change is prepared before it happens \cite{riverso2013plug, zeilinger2013plug,
aboudonia2022reconfigurable, saccani2023model}, all on linear coupled subsystems or ground robots; a suspended payload
couples every vehicle through the load, and a magnet that lets go happens at an instant the planner does not choose, so
this thesis switches to a pre-built problem of the new size at the event and shows the first solve in flight rather
than certifying it. Static equilibrium decides which layouts can hold a payload at all: three robots are the fewest,
their attachment points must span its centre of mass \cite{michael2011cooperative, jiang2013inverse}, and the shares
are exact when the centre lies in the attachment plane \cite{tognon2018internal}; see also
\cite{gabellieri2023equilibria, liu2025wrench, fink2011planning}. All analyse a fixed team; Section~\ref{sec:reconfig}
turns the condition into a margin that says which carrier a layout may lose.

\begin{table}[htbp]
\centering
\footnotesize
\setlength{\tabcolsep}{3pt}
\caption[Changes of team membership under a carried payload]{Changes of team membership under a carried
payload, as each work reports them. Announced: the controller is told before it happens. Leaver: whether the vehicle
leaves the load or stays attached. Detection: from the event to the controller acting on it; recovery: to a settled
state. ``--'': not reported; brackets: to fill from the full texts.}
\label{tab:membership}
\begin{tabularx}{\textwidth}{@{}
  >{\raggedright\arraybackslash\hsize=0.86\hsize}X
  >{\raggedright\arraybackslash\hsize=1.04\hsize}X
  >{\raggedright\arraybackslash\hsize=0.78\hsize}X
  >{\raggedright\arraybackslash\hsize=0.90\hsize}X
  >{\raggedright\arraybackslash\hsize=1.04\hsize}X
  >{\raggedright\arraybackslash\hsize=1.14\hsize}X
  >{\raggedright\arraybackslash\hsize=1.24\hsize}X@{}}
\toprule
Work & Event; mechanism & Announced & Leaver & Plant; team; payload & Attitude or sag through the event &
Detection; recovery \\
\midrule
Quan\newline \cite{quan2025swarm}, 2025 & release on command; cable released at the load & yes & removed, flies off with
its cable & hardware, outdoor; 4 to 3; 4\,kg & -- (tensions and one robot's velocity plotted) & --; tensions
balanced after 13\,s \\
\addlinespace
Zeng\newline \cite{zeng2025marl}, 2025 & motors switched off & no & left hanging from the load & hardware; 3 to 2;
1.4\,kg & plotted, no number & none explicit (learned policy); -- \\
\addlinespace
Bacheti\newline \cite{tensionbased}, 2023 & planned release; the team repositions first & yes & [removed; count to
verify] & hardware; [3 to 2 to 3]; point mass & -- & --; -- \\
\addlinespace
Liang\newline \cite{liang2021fault}, 2021 & cable breaks & no & removed & simulation; [team]; [mass] & [from the paper] &
disturbance observers; [from the paper] \\
\addlinespace
Saiella\newline \cite{saiella2021fault}, 2021 & vehicle fault; the supervisor reconfigures & no & [from the paper] &
simulation; [team] & -- & supervisor; -- \\
\addlinespace
Hajieghrary\newline \cite{hajieghrary2026passive}, 2026 (preprint) & one or two cables severed & no & removed &
simulation; 5 to 4 and 5 to 3 & [from the preprint] & none (passive feed-forward); within a pendulum period
[0.32\,s] \\
\addlinespace
Farras\newline \cite{farrasFaultTolerantControlMultiQuadcopter2024}, 2024 & [fault, from the paper] & no & [from the
paper] & simulation; [team] & [from the paper] & [from the paper] \\
\addlinespace
Shankar\newline \cite{shankar2020inflight}, 2020 & a bar handed from one multirotor to another & yes & the giver releases
after the grasp is confirmed & hardware; 1 to 1 & -- & magnetic sensor at the grasp; gains rescheduled \\
\addlinespace
ModQuad\newline \cite{saldana2018modquad}, TRADY \cite{sugihara2023trady}, LEGION \cite{sugihara2026legion}, Miyazaki \cite{miyazaki2018airborne} & vehicles docking to one another; no payload & yes & n/a & hardware; 1 to 1 per dock &
TRADY: altitude jump 0.39\,m, 0.065\,m with the thrust blend & motion capture, contact force or a latch; -- \\
\midrule
This work & magnet off, commanded or not (leave); contact and co-motion (join) & both & removed; lands by itself
& hardware and twin; 4 to 3, 3 to 4 (twin); 0.86\,kg ring on rods & tilt rise 5.3--16.9$^\circ$ over five rig
releases & 0.33--0.42\,s (rig); join: [join tree] \\
\bottomrule
\end{tabularx}
\end{table}


% lit_capture.tex -- Literature 2.3 (v3_50: condensed).

\subsection{Capture proof, docking and magnetic grasping}
\label{sec:lit_capture}

A join needs a mechanism that catches, a proof that it has caught, and a switch of control at the proof. A grasp can
be proven by a sensor in the gripper \cite{shankar2020inflight}, or from the vehicles' motion: ModQuad
\cite{saldana2018modquad} detects the new structure from motion capture once the docked modules move as one body, and
LEGION \cite{sugihara2026legion} treats a latched joint as proven when the contact wrench settles. A proof from pose is
the only kind available here, since nothing on the ring or the magnet is instrumented (Section~\ref{sec:reconfig}).
Modular vehicles dock to one another in flight, often through magnets, and the controller switch matters as much as
the mechanism: TRADY \cite{sugihara2023trady} cut the altitude jump after assembly sixfold by blending thrust from the old
controller to the new one; see also \cite{zhang2026epm, miyazaki2018airborne, antal2026hook, xu2026hitches}. Each
docks to a rigid structure or a single carrier; none joins a payload hung from a team. A magnet holds best when pulled
straight off its plate and loses hold as the pull turns across the face \cite{gawel2017magnetic}; the rig's magnets catch
on contact only and every ring lost on the rig began with a magnet letting go (Section~\ref{sec:res_leave}), which is why
the rods are planned steep, announced changes are ramped and a newcomer is folded in from a small tension.


% lit_switching.tex -- Literature 2.4 (v3_50: condensed).

\subsection{Switching, bumpless transfer and runtime assurance}
\label{sec:lit_switching}

A membership change is a switch between two plants and so between two controllers. When control passes from one to
another, the second must start from the first's output or the plant takes a step; the classical remedy keeps the idle
controller tracking the active one \cite{graebe1996transfer, peng1996antiwindup}. Here the new problem is warm-started
from the measured state with the tensions free at its first node, an announced change is ramped, and a newcomer is
folded in from a small tension (Section~\ref{sec:reconfig}). Nonlinear MPC in flight runs the real-time iteration of
Diehl et al.\ \cite{diehl2002rti, verschueren2022acados}, whose accuracy rests on the previous solution being close to
the next; a switch to a problem of a different size breaks that chain \cite{gros2020rti}, so each solver is built in
advance and primed before it flies, and the first live solve after a switch is reported as the switch's validation
(Section~\ref{sec:res_switch}). A complex controller can be flown inside a simple safety envelope: the Simplex
architecture hands control to a verified baseline when a monitor trips \cite{sha2001simplicity}, and the fall-backs of
Section~\ref{sec:safety} are framed the same way. A system whose agents join and leave while it runs is an open
multi-agent system \cite{hendrickx2017open, sekercioglu2026open}; on a carried payload the agents are coupled through
one rigid body, so a change of membership is also a change of plant.


% lit_execution.tex -- Literature 2.5 (v3_50: condensed).

\subsection{Execution layer and force estimation}
\label{sec:lit_execution}

In cooperative transport the execution layer between a planner's references and each vehicle's motors has one job
the single-vehicle case does not: the link's pull. In Sun et al.\ \cite{tudelft} an onboard INDI controller
\cite{smeur2016indi} compensates the cable force from the vehicle's own accelerometer; estimating the force is the other
route \cite{klausen2020cooperative, tagliabue2019robust}, and shared loads have been flown with the execution onboard
each rotorcraft \cite{demarina2019flexible, bulka2022experiments}. Each needs code on the vehicle. The vehicles of this
thesis run unmodified Betaflight, which closes a body-rate loop and accepts a throttle and three rate commands over a
radio link, so the execution layer is a tracking MPC per vehicle on the ground, with the rod's pull entering its model
as a known input: from the planner's tensions, or measured from the vehicle's motion-capture acceleration minus its
modelled thrust (Section~\ref{sec:tracker}). What the measured pull buys is the first research question.

\paragraph{Steady-error rejection.} Offset-free MPC removes the steady error of a constant disturbance by adding it to
the prediction model and estimating it from the measurements \cite{muske2002disturbance, pannocchia2003disturbance,
maeder2009linear}; by the internal model principle \cite{francis1976internal} a constant is then rejected, and whether
it is modelled at the input or the output changes the transient, not the steady state \cite{pannocchia2007combined}.
Section~\ref{sec:disturbance} builds the planner's integral action on this.


% lit_sim.tex -- Literature 2.6 (v3_50: condensed).

\subsection{Simulation to reality}
\label{sec:lit_sim}

The most faithful simulation reviewed runs the platform's own estimator, controllers and motor firmware inside Gazebo
\cite{distributedNMPC, koenig2004gazebo}; linearised rigid-rod models allow rapid iteration but miss the large attitudes
of a transient \cite{DistributedMPC}. Simulators built for suspended loads report their fidelity per axis against
hardware \cite{li2024rotortm, furrer2016rotors, cruz2017cable}. This thesis uses a Gazebo twin of the rig built from the
rig's identified thrust map, geometry and masses, with an emulated flight controller, and adds two things to the usual
practice: the twin's omissions are stated before the results, and the twin is to be scored as a predictor, each rig cell
against a band frozen before it flies (Section~\ref{sec:ledger}).


\subsection{The gap}
\label{sec:lit_gap}

Three things are missing from the literature above, and the three research questions of Section~\ref{sec:rq}
follow from them. No work attaches a vehicle to a payload that a team is already carrying, in simulation or on
hardware: the docking literature joins vehicles to vehicles or to a structure, and the transport literature
describes a joining vehicle in prose at most. No hardware work read in full reports the peak attitude transient of
a true release from a team-carried payload, and where a time is given it is the time to recover, not the time to
detect; the simulation work on severed cables reports recovery too, and none of it flies. And no work scores its
simulator as a predictor of its own hardware flights, with the prediction written down before the flight: the
usual practice is to report agreement afterwards. This thesis flies the leave on hardware and measures the ring's
attitude through it, builds and tests an admission contract for a vehicle flown by another controller, and freezes
the twin's prediction for each rig flight before it flies.
